\subsection{Comparison with the known eleven-vertex Landau--Lifshitz pair}
\label{sec:class6-eleven-vertex-comparison}

The Class 6 lattice model is quantum-mechanically the eleven-vertex model
\cite{CorcoranDeLeeuw2024}, and the associated rational Landau--Lifshitz
hierarchy is already known \cite{LevinOlshanetskyZotov2014}.  We now compare
the pair derived above directly with an explicit form of that known classical
pair.  This provides an independent check of the finite-lattice construction;
it is not a claim of a new Class 6 hierarchy.

Following the rational formulas of Ref.~\cite{AtalikovZotov2021}, introduce
a traceless spin matrix
\begin{equation}
\Sigma=
\begin{pmatrix}
\Sigma_{11}&\Sigma_{12}\\
\Sigma_{21}&-\Sigma_{11}
\end{pmatrix}
\end{equation}
and the spectral parameter $z$.  The spatial matrix can be written as
\begin{equation}
U_{\rm 11v}(z,\Sigma)=
\begin{pmatrix}
\Sigma_{11}/z-z\Sigma_{12}&\Sigma_{12}/z\\
\Sigma_{21}/z-2z\Sigma_{11}-z^3\Sigma_{12}
&-\Sigma_{11}/z+z\Sigma_{12}
\end{pmatrix}.
\label{eq:class6-known-eleven-U}
\end{equation}
For $\Sigma^2=\Lambda_{\Sigma}^2\Id_2$, define
\begin{align}
M_{\rm top}(z,\Sigma)
&=\begin{pmatrix}
\Sigma_{12}&0\\
2\Sigma_{11}+2z^2\Sigma_{12}&-\Sigma_{12}
\end{pmatrix},
\nonumber\\
\mathsf v
&=-\frac{k}{4\Lambda_{\Sigma}^2}
[\Sigma,\partial_x\Sigma].
\end{align}
The corresponding time matrix is
\begin{equation}
V_{\rm 11v}(z,\Sigma)
=\frac12\left[
\frac{1}{z}U_{\rm 11v}(z,\Sigma)
+2M_{\rm top}(z,\Sigma)
-U_{\rm 11v}(z,\mathsf v)
\right],
\label{eq:class6-known-eleven-V}
\end{equation}
with zero-curvature convention
$\partial_{\tau}U_{\rm 11v}+k\partial_xV_{\rm 11v}
-[U_{\rm 11v},V_{\rm 11v}]=0$.

For the Class 6 spin field, set
\begin{equation}
\Sigma:=
\begin{pmatrix}
S^3&S^+\\
S^-&-S^3
\end{pmatrix}.
\label{eq:class6-eleven-spin-matrix}
\end{equation}
The unit-spin constraint gives $\Sigma^2=\Id_2$.  Let $\beta$ be a nonzero
complex number satisfying
\begin{equation}
\beta^2=-2\alpha_6,
\qquad
z=\beta\lambda,
\qquad
k=-\frac{4}{\beta},
\qquad
\tau=\frac{\ii\beta^2}{2}t=-\ii\alpha_6t.
\label{eq:class6-eleven-dictionary-paper}
\end{equation}
Then the traceless spatial matrix $U_{6,0}$ defined in
Eq.~\eqref{eq:class6-traceless-U} and the time matrix
Eq.~\eqref{eq:class6-time-Lax-compact} obey the exact identities
\begin{align}
U_{6,0}(\lambda)
&=\frac{\beta}{4}
U_{\rm 11v}(\beta\lambda,\Sigma)^{\mathsf T},
\label{eq:class6-eleven-U-map-paper}\\
V_6(\lambda)
&=\frac{\ii\beta^2}{2}
V_{\rm 11v}(\beta\lambda,\Sigma)^{\mathsf T}.
\label{eq:class6-eleven-V-map-paper}
\end{align}
The scalar difference $U_6-U_{6,0}=\Id_2/(4\lambda)$ does not affect the
curvature.  The transpose reverses the commutator order, while
$(\beta/4)k=-1$ and the time rescaling in
Eq.~\eqref{eq:class6-eleven-dictionary-paper} reproduce the zero-curvature
convention used in this paper.

The same dictionary maps the known eleven-vertex spin equation to
Eq.~\eqref{eq:class6-vector-eom}.  Thus the spatial matrix, the time matrix
and the Hamiltonian flow obtained from the Class 6 quantum chain all agree
with the independently known rational eleven-vertex Landau--Lifshitz
structure.
